\chapter{Theoretical Background}
\label{ch:theory}
\section{The beauty of Symmetry}

Over the centuries, physics has developed increasingly with the aim of understanding and describing the universe and its nature. One might think that mathematical complexity has increased steadily in order to describe the universe, and in a sense this is true, but how far have we gone? Is there a limit to the complexity we can achieve? Is there a risk, however remote, that physicists might lose control of this mathematical machine, which is becoming ever more complex and detached from any conceivable reality and, ultimately, from the physical world?

\noindent To answer this question, I would like to address the topic in physics that has most captured the author's attention: \textit{the role of symmetry.}


\subsubsection{Simplicity Behind Complexity}
\begin{center}
    `The bigger the puzzle, the harder it is to solve.'
\end{center}

There is certainly some truth in this statement. To put it into context, let's consider an example. If I wanted to describe a room, I would have to refer to every object inside it, and then specify its position and orientation. In short, I could spend hours describing a room in minute detail. If I then wanted to describe the entire flat, then move on to the building, and then go on to describe a city, the task would become increasingly complex. At this point, the more attentive readers will already have guessed where I'm heading.

Describing the universe seems an impossible task, unless we find a way to simply represent something far greater than anything simple. This is where physicists' interest in symmetries stems from.

Let's return to a concrete example: a beehive. I could make the task of describing it very long and tedious, scrutinising every single detail that makes it up. Or I could note a pattern that extends throughout the entire beehive. All I need to do is describe a single hexagon and then observe that, by symmetry, the rest of the beehive follows the same structure.

But what exactly is symmetry? We can define it in various ways; let's look at two of them. From an intuitive point of view, symmetry is a pattern that exhibits a certain regularity in space and time. From a physical point of view, we can define it as a property that remains unchanged following a transformation.

Thus begins the `hunt for symmetries' by physicists. I would like to mention some of the most common ones in physics. If we think of a wheel, it possesses rotational symmetry about an axis, the one passing through its hub. Another fundamental symmetry is temporal symmetry: if we shift a system forwards or backwards in time, the laws of physics must remain the same.  Without this symmetry, physics would be unable to make predictions and therefore formulate universal laws.

Nature possesses many other symmetries. There is, for example, reflection symmetry: like in a mirror, it allows us to recognise when a configuration remains equivalent to its reflected image.

In the next section, we will delve into the mathematics that allows us to describe this fundamental concept, which underpins the evolution of physics and, with it, our understanding of the universe.

\begin{comment}{
"The bigger the puzzle, the harder to solve."
 Asserzione che sicuramente ha del vero. Per contestualizzarla, consideriamo un esempio. Se volessi descrivere una stanza dovrei far riferimento ad ogni oggetto al suo interno, per poi specificarne la posizione e l'orientamento. Insomma potrei passare ore a descrivere minuziosamente una stanza. Se poi volessi descrivere l intero appartamento, per poi passare alla palazzina, per poi passare alla descrizione di una città, il compito diventerebbe sempre più complesso. A questo punto, i lettori più attenti avranno già intuito dove voglio arrivare.

Descrivere l'universo sembra un impresa irrangiungibile, a meno che non troviamo un escamotage per rappresentare con semplicità qualcosa di assai più grande che di semplice non ha nulla. Da qui nasce l'interesse dei fisici per le simmetrie.

Torniamo ad un esempio concreto, ovvero un alveare. Potrei rendere il compito di descriverlo molto lungo e tedioso, guardando miniziosamente ad ogni dettaglio che lo compone. Oppure potrei notare un pattern che si estende per tutto l'alveare. Mi basta descrivere un singolo esagono e poi osservare che, per simmetria, il resto dell'alveare segue la stessa struttura.

ma cosa è esattamente una simmetria? possiamo definirla in vari modi, vediamone due. Dal punto di vista intuitivo, una simmetria è un pattern che presenta una certa riccorenza nello spazio e nel tempo. Dal punto di vista fisico, possiamo definirla come una proprietà che rimane invariata in seguito a una trasformazione.

Inizia cosi la "caccia alle simmetrie" da parte dei fisici. Vorrei citarne alcune tra le più ricorrenti in fisica. Se pensiamo ad una ruota, essa possiede una simmetria di rotazione attorno ad un asse, quello che passa attraverso il suo perno. Un altra simmetria fondamentale è quella temporale, se trasliamo un sistema avanti o indietro nel tempo, le leggi fisiche devono rimanere le stesse.  Senza questa simmetria, la fisica non sarebbe in grado di fare predizioni e quindi formulare leggi universali.

La natura possiede molte altre simmetrie. C'è, per esempio, la simmetria di riflessione: come in uno specchio, essa permette di riconoscere quando una configurazione rimane equivalente alla propria immagine riflessa.

Nella prossima sezione ci addentreremo nella matematica che permette di descrivere questo concetto fondamentale, che tesse l'evoluzione della fisica e, con essa, della nostra comprensione dell'universo.}
\end{comment}

\section{Group Theory}
\label{sec:theory-group}

The previous section argued that symmetry is what makes a description of the
universe tractable. This section makes that statement precise, and it does so
with a very concrete goal in mind. Every field in the models of this thesis
carries indices, and the toolkit of \Cref{ch:feynpy} asks the user to declare
each of them by type. Those types are not arbitrary labels: a colour index, a
weak-isospin index and an adjoint index are all representation indices of an
underlying Lie algebra, and it is the representation that decides which
indices may be contracted with which. The path from here to there runs
through Lie groups, Lie algebras, generators and structure constants, and
finally representations.

Symmetries are ubiquitous in physics, and since physics describes nature
through mathematics, symmetries too are described mathematically, through the
concept of a group. They may be discrete, such as the lattice symmetries of
condensed matter, or --- more interesting for our purposes --- continuous. In
the latter case they are known as Lie groups. Lie groups play a double role in
quantum field theory: they appear as global spacetime and flavour symmetries,
and they appear as local gauge symmetries in Yang--Mills theory, which is the
case that will occupy us for the rest of the thesis.

\subsection{Lie Groups and Lie algebras}
 we start by defining Lie groups, Lie algebras and their relationship.

 \begin{definition}[Lie group]
A Lie group $G$ is a group that is also a smooth manifold, such that
the group multiplication $G \times G \to G$ is a smooth map.
\end{definition}
Lie groups are curved manifolds which makes them hard to study, therefore we look at the tangent space at the identity which, equipped with a natural product is called Lie algebra.
\begin{definition}[Lie algebra]
The Lie algebra $\mathfrak{g}$ associated to a Lie group $G$ is the
tangent space at the identity element $1 \in G$:
\begin{equation}
    \mathfrak{g} := T_1 G.
\end{equation}
\end{definition}


Elements of $\mathfrak g$ should be thought of as
infinitesimal group elements, in a more physical wording they are generators of infinitesimal
transformations.

Having passed from the Lie group $G$ to its Lie algebra
$\mathfrak{g} = T_1 G$, we now ask whether this can be reversed: is
there a map that takes us back from $\mathfrak{g}$ to $G$, at least in
a neighbourhood of the identity? Such a map indeed exists and is
called the exponential map.

\begin{definition}[Exponential map]
The exponential map is the smooth map from a neighbourhood of
$0 \in \mathfrak g$ to a neighbourhood of $1 \in G$ satisfying
\begin{equation}
    \exp(0) = 1, \qquad d\exp(0) = \mathrm{id}, \qquad \exp(na) = \exp(a)^n.
\end{equation}
\end{definition}

It can be constructed explicitly as the limit
\begin{equation}
    \exp(a) = \lim_{n\to\infty}\left(1+\frac{a}{n}\right)^n.
\end{equation}

\begin{definition}[Lie bracket]
The Lie bracket $\llbracket \cdot,\cdot \rrbracket$ is (twice) the
leading deviation of $m$ from linearity, $m_2(a,b) =
\tfrac12 \llbracket a,b \rrbracket$. It can equivalently be recovered
from a group commutator of exponentials,
\begin{equation}
    \llbracket a,b \rrbracket
    = \lim_{\epsilon \to 0} \frac{1}{\epsilon^2}
      \exp^{-1}\!\Big(\exp(\epsilon a)\exp(\epsilon b)\exp(-\epsilon a)\exp(-\epsilon b)\Big).
\end{equation}
\end{definition}

The bracket therefore directly measures the failure of $G$ to be
commutative; an abelian group has vanishing bracket.

The Lie bracket satisfies three key properties, each of which will be
used repeatedly in what follows. It is bilinear, and it is
antisymmetric,
\begin{equation}
    \llbracket a,b \rrbracket = -\llbracket b,a \rrbracket,
\end{equation}
and it satisfies the Jacobi identity,
\begin{equation}
    \big\llbracket \llbracket a,b\rrbracket, c\big\rrbracket
    + \big\llbracket \llbracket b,c\rrbracket, a\big\rrbracket
    + \big\llbracket \llbracket c,a\rrbracket, b\big\rrbracket = 0.
\end{equation}
The Jacobi identity is the algebraic shadow of associativity of the
group multiplication; it is not an extra assumption but a consequence
of $G$ being a group.

\begin{definition}[Lie algebra]
A vector space $\mathfrak g$ equipped with a bracket satisfying
bilinearity, antisymmetry and the Jacobi identity is called a Lie
algebra. Conversely, the exponential map associates to every such
$\mathfrak g$ a (simply connected) Lie group $G$.
\end{definition}

From here on we work directly with the algebraic data
$(\mathfrak g, \llbracket \cdot,\cdot\rrbracket)$ rather than with the
group manifold itself.

% --- 4. Basis, generators, structure constants ---------------------
\subsection{Basis, generators, structure constants}

Now we introduce a key concept that will allow us to understand how
indices actually work. If we decompose everything with a chosen basis $\{T_a\}$ of $\mathfrak g$, we get some useful notation, in particular for the lie bracket. The elements $\{T_a\}$ are called generators of of $\mathfrak g$. In physics the basis is typically chosen to be imaginary
\begin{equation}
    a = i\, a^c T_c \in \mathfrak{g},
\end{equation}
This convention makes $T_a$ hermitian which is nice because it guarantees real eigenvalues. The downside is that we get many factors of $i$. For instance, a group element is parametrised as
\begin{equation}
    h = \exp\!\big(i\, a^c T_c\big) \in G.
\end{equation}

\begin{definition}[Structure constants]
Expanding the Lie bracket of two generators in the chosen basis
defines the structure constants $f_{ab}{}^{c}$ of the algebra,
\begin{equation}
    \llbracket T_a, T_b \rrbracket = i f_{ab}{}^{c} \, T_c.
\end{equation}
\end{definition}

For a real Lie algebra $f_{ab}{}^{c}$ must be real numbers, and
antisymmetry of the bracket implies
\begin{equation}
    f_{ab}{}^{c} = -f_{ba}{}^{c}.
\end{equation}
This index $a$, labelling a generator $T_a$, is what we will call an
index or, once specialised to $SU(N)$ gauge theory, a colour index.
Once we pass to representations, these same structure constants will
reappear as the commutation relations of the representation matrices
$T_a^R$.

% --- 5. Representations: why fields carry indices -------------------
\subsection{Representation}

In Yang--Mills theory fields have particular transformation
properties under gauge transformations. Mathematically these
transformation rules are described by representations.

\begin{definition}[Representation of a Lie group]
A representation $R$ of $G$ is a map from the group to automorphisms
of some vector space $V$,
\begin{equation}
    R : G \to \mathrm{Aut}(V).
\end{equation}
\end{definition}

\begin{definition}[Representation of a Lie algebra]
A representation $R$ of $\mathfrak g$ is a linear map from the
algebra to endomorphisms of some vector space $V$,
\begin{equation}
    R : \mathfrak{g} \to \mathrm{End}(V),
\end{equation}
such that the Lie bracket is reproduced as an operator commutator,
\begin{equation}
    R\big(\llbracket a,b \rrbracket\big) = [R(a), R(b)].
\end{equation}
\end{definition}

Writing $T_a^R := R(T_a)$ for the representation matrices of the
generators, the representation property of $R$ takes the explicit form
\begin{equation}
[T_a^R, T_b^R] = i f_{ab}{}^{c}, T_c^R .
\end{equation}
A field $\phi$ transforming in the representation $R$ carries an
index $i = 1, \dots, \dim R$ labelling a basis of the representation
space $V$, and an infinitesimal gauge transformation acts on it as
\begin{equation}
\delta \phi^i = i a^c (T_c^R)^i{}_j \phi^j .
\end{equation}
This is the sense in which internal indices in quantum field theory,
such as colour or weak-isospin indices, are representation indices of
the underlying symmetry algebra.

Now we want to study a special representation: the adjoint representation. Its feature is that the Lie algebra acts on itself.

\begin{definition}[Adjoint representation]
A distinguished representation of any Lie algebra $\mathfrak g$ is the adjoint representation on the Lie algebra itself:
\begin{equation}
\operatorname{ad}: \mathfrak g \to \operatorname{End}(\mathfrak g),
\qquad
\operatorname{ad}(a)b \coloneqq [a,b].
\end{equation}
\end{definition}

In other words, an element of the algebra acts on another element by taking the commutator. If ${T_a}$ is a basis of the Lie algebra, then
\begin{equation}
T^{ad}_aT_b=\operatorname{ad}(T_a)T_b = \llbracket T_a,T_b\rrbracket = i f_{ab}{}^cT_c.
\end{equation}
Therefore, the adjoint representation $T^{ad}_a$ is a matrix whose elements are the structure constants

\begin{equation}
(T^{ad}_a)_b{}^c = i f_{ab}{}^{c},
\end{equation}


Thus, an adjoint index is an index labelling a basis element of the Lie algebra. For example, the gauge field can be written as
\begin{equation}
A_\mu = A_\mu^a T_a .
\end{equation}
Here $\mu$ is a Lorentz index, while $a$ is an adjoint index. The index $a$ labels which generator of the gauge algebra the field component is associated with.


\subsection{Invariants}

Lie algebras possess an invariant object which lets us raise and
lower algebra indices, in the same way that a metric lets us raise
and lower spacetime indices.

\begin{definition}[Killing form]
A simple Lie algebra $\mathfrak g$ has a unique invariant symmetric
bilinear form $K : \mathfrak g \times \mathfrak g \to K$, the Killing
form, satisfying
\begin{equation}\label{killing_prop}
    K(a,b) = K(b,a), \qquad
    K\big(\llbracket c,a\rrbracket, b\big) + K\big(a, \llbracket c,b\rrbracket\big) = 0
\end{equation}
for all $a,b,c \in \mathfrak g$.
\end{definition}
A form with this property can be constructed easily using some representation $R$
\begin{equation}
    Tr R(a)R(b) =B^R K(a,b)
\end{equation}
The property in equation \ref{killing_prop} follow from cyclicity of the trace. Due to uniqueness of
$K$ in a simple Lie algebra, all these forms must be equivalent up to a factor of proportionality $B^R$ which depends on the particular representation R.
In the basis of generators, the Killing form has components
\begin{equation}
    K(T_a, T_b) = k_{ab}.
\end{equation}
For a real, compact, simple Lie algebra $k_{ab}$ is positive-definite
and non-degenerate, and can therefore be used to raise and lower the
algebra index $a$. In particular we may lower an index on the
structure constants,
\begin{equation}
    f_{abc} := f_{ab}{}^{d}\, k_{dc},
\end{equation}
which are then totally antisymmetric in all three indices,
\begin{equation}
    f_{abc} = -f_{bac} = -f_{acb}.
\end{equation}
A common simplifying choice of basis sets $k_{ab} = \delta_{ab}$, so
that upper and lower algebra indices coincide numerically.

\subsection{The Unitary Algebra \texorpdfstring{$\mathfrak{su}(N)$}{su(N)}}

We specialise the index calculus developed above to the case of
$\mathfrak{su}(N)$, the gauge algebra of quantum chromodynamics for
$N = 3$.

\begin{definition}[Defining representation]
$\mathfrak{su}(N)$ acts naturally on $\mathbb{C}^N$ by matrix
multiplication,
\begin{equation}
    \mathrm{def}(A)\, v = Av, \qquad v \in \mathbb{C}^N, \qquad D^{\mathrm{def}} = N.
\end{equation}
\end{definition}

This is the representation carried by matter fields such as quarks;
the corresponding index, written $i = 1,\dots,N$, is the colour
index. It is a different kind of index from the adjoint index $a$
carried by the gauge field, and the two are related by the
generators $T_a^{\mathrm{def}}$, which carry one of each,
$(T_a^{\mathrm{def}})^i{}_j$.

Together with the identity matrix, the generators
$T_a^{\mathrm{def}}$ form a complete basis of $N \times N$ matrices.
This completeness is expressed by the relation
\begin{equation}
    k^{ab} (T_a^{\mathrm{def}})^i{}_j (T_b^{\mathrm{def}})^k{}_l
    = B^{\mathrm{def}} \left(\delta^i_l \delta^k_j - \frac{1}{N}\delta^i_j \delta^k_l \right),
\end{equation}
and it is this identity that allows a sum over the adjoint index $a$
to be traded for an expression written purely in terms of Kronecker
deltas on the colour indices $i,j,k,l$ --- a manipulation used
throughout when working with $SU(N)$ gauge theories.


% ==============================================================
\section{Gauge Theories and Their Lagrangians}
\label{sec:theory-gauge}
% ==============================================================

We now have the algebraic vocabulary we need. This section assembles it into
the object the rest of the thesis manipulates: the Lagrangian of a gauge
theory. The presentation is deliberately structural rather than
pedagogical --- the aim is to fix which pieces exist, what indices they
carry, and how they are built, because those are exactly the pieces that
\Cref{ch:feynpy} asks the user to declare and \Cref{ch:internals} teaches the
compiler to expand. Standard textbook treatments are
in~\cite{Peskin:1995ev,Weinberg:1996kr}.

\subsection{Fields and the indices they carry}
\label{sec:theory-fields}

A field is characterised by how it transforms, and every transformation law
it obeys leaves behind an index. There are three families.

\emph{Spacetime} transformations fix the spin. A scalar $\phi$ carries no
spacetime index; a Dirac fermion $\psi_\alpha$ carries a four-component
spinor index $\alpha$; a vector field $A_\mu$ carries a Lorentz index $\mu$.
The Dirac matrices $\gamma^\mu_{\alpha\beta}$ carry one of each and are
therefore the object that ties spinor lines to Lorentz lines.

\emph{Internal gauge} transformations fix the representation, in exactly the
sense of \Cref{sec:theory-group}. A quark $q^i$ carries a colour index in the
fundamental of $\mathrm{SU}(3)_c$; a gluon $G^a_\mu$ carries an adjoint index;
a left-handed lepton doublet $\ell_L^j$ carries a weak-isospin index in the
fundamental of $\mathrm{SU}(2)_L$. An abelian group such as
$\mathrm{U}(1)_Y$ leaves no index at all, only a number: the charge.

\emph{Flavour} indices label copies of an otherwise identical field. The
three generations of quarks and leptons are the standard example, and the
index they carry is contracted by coupling matrices such as the Yukawa
matrices $Y_u, Y_d, Y_e$ rather than by a gauge generator.

A single field can carry one index of each family at once. The left-handed
quark doublet has a spinor index, a colour index, a weak index and a flavour
index simultaneously, and it is precisely this multiplicity that makes
hand-derivation of Feynman rules so unpleasant. The toolkit's
\texttt{IndexType}, introduced in \Cref{sec:feynpy-indices}, is nothing more
than the machine-readable version of the classification just given.

Chirality deserves a separate remark because it recurs throughout the
validation of \Cref{sec:validation-chirality}. The projectors
\begin{equation}
    P_L = \frac{1-\gamma^5}{2}, \qquad P_R = \frac{1+\gamma^5}{2},
    \label{eq:theory-projectors}
\end{equation}
split a Dirac field into left- and right-handed parts. In the Standard Model
the two chiralities transform differently under $\mathrm{SU}(2)_L$, so the
natural fields of the unbroken theory are chiral from the start: the doublets
$Q_L$ and $\ell_L$ and the singlets $u_R$, $d_R$, $e_R$. A theory written in
terms of such fields carries its chirality information in the field names,
and an explicitly written projector then becomes redundant --- a fact the
comparison of \Cref{ch:validation} both exploits and verifies.

\subsection{Local symmetry and the covariant derivative}
\label{sec:theory-covd}

Let a field $\phi$ transform in a representation $R$ of a group $G$, so that
under a global transformation
\begin{equation}
    \phi \longrightarrow e^{i \alpha^a T_a^R}\, \phi ,
\end{equation}
with constant parameters $\alpha^a$. A Lagrangian built from $\phi$ and
$\partial_\mu \phi$ is invariant because the derivative passes through the
constant matrix untouched.

Promoting the symmetry to a \emph{local} one, $\alpha^a \to \alpha^a(x)$,
breaks that argument: the derivative now also hits the parameter, and
$\partial_\mu \phi$ no longer transforms in the representation $R$. The
repair is to introduce a gauge field $A^a_\mu$, one component per generator,
and replace the ordinary derivative by the covariant derivative
\begin{equation}
    D_\mu \phi = \partial_\mu \phi - i g\, A^a_\mu\, T_a^R\, \phi ,
    \label{eq:theory-covd}
\end{equation}
where $g$ is the coupling constant of the group. If the gauge field is
declared to transform inhomogeneously, its variation cancels the unwanted
term and $D_\mu \phi$ transforms exactly like $\phi$ itself. For an abelian
group there is a single generator and \Cref{eq:theory-covd} reduces to
$D_\mu\phi = \partial_\mu\phi - i g Q A_\mu \phi$ with $Q$ the charge of the
field.

Two features of \Cref{eq:theory-covd} matter for the implementation. First,
the sign convention is a choice, and it must be made once and enforced
everywhere; the convention used throughout this thesis is the one written
above. Second, a field charged under several groups picks up one term per
group, each with its own coupling, its own gauge field, and its own generator
acting on its own index. The left-handed quark doublet gives
\begin{equation}
    D_\mu Q_L
    = \partial_\mu Q_L
    - i g_s\, G^a_\mu\, t^a_{(3)}\, Q_L
    - i g\, W^I_\mu\, t^I_{(2)}\, Q_L
    - i g'\, Y_Q\, B_\mu\, Q_L ,
    \label{eq:theory-covd-sm}
\end{equation}
a single symbol standing for four separate terms. Expanding that symbol
correctly, for any field and any set of groups, is the job of
\Cref{sec:internals-covd}.

\subsection{The field strength and the Yang--Mills Lagrangian}
\label{sec:theory-fieldstrength}

The gauge field needs a kinetic term of its own, and it must be built from a
gauge-covariant object. That object is the field strength, defined by the
commutator of two covariant derivatives,
\begin{equation}
    F_{\mu\nu} = \frac{i}{g}\,[D_\mu, D_\nu] = F^a_{\mu\nu} T_a ,
\end{equation}
whose components read
\begin{equation}
    F^a_{\mu\nu}
    = \partial_\mu A^a_\nu - \partial_\nu A^a_\mu
    + g\, f^{abc}\, A^b_\mu A^c_\nu .
    \label{eq:theory-fs}
\end{equation}
The last term is present only for a non-abelian group, and it is the origin
of the gauge boson self-interactions: the Yang--Mills Lagrangian
\begin{equation}
    \mathcal{L}_{\mathrm{YM}} = -\frac14 F^a_{\mu\nu} F^{a\,\mu\nu}
    \label{eq:theory-ym}
\end{equation}
is quadratic in $F$, and therefore contains a two-gluon kinetic piece, a
three-gluon vertex and a four-gluon vertex all at once. For an abelian group
the structure constants vanish, \Cref{eq:theory-fs} reduces to
$F_{\mu\nu}=\partial_\mu A_\nu - \partial_\nu A_\mu$, and the photon does not
interact with itself.

The counting implicit in \Cref{eq:theory-fs} recurs in
\Cref{sec:internals-fs}: a non-abelian field strength is a sum of three
pieces, so a product of $n$ of them expands into $3^n$ branches before
anything is collected again.

\subsection{Matter, Yukawa couplings and flavour}
\label{sec:theory-matter}

With covariant derivatives available, matter Lagrangians are short. A Dirac
fermion contributes
\begin{equation}
    \mathcal{L}_\psi = i\,\bar\psi\,\gamma^\mu D_\mu \psi ,
    \qquad
    \bar\psi = \psi^\dagger \gamma^0 ,
\end{equation}
and a complex scalar contributes $(D_\mu\Phi)^\dagger (D^\mu\Phi)$. Expanding
$D_\mu$ in either expression produces the kinetic term plus one gauge current
per group; the interaction between matter and gauge fields is entirely
contained in that expansion and is never written separately.

Fermion masses cannot be written directly in a chiral gauge theory, because
$\bar\psi_L \psi_R$ is not gauge invariant when the two chiralities sit in
different representations. They arise instead from Yukawa couplings to a
scalar,
\begin{equation}
    \mathcal{L}_Y = - (Y_u)_{fg}\, \bar Q_L^f\, \tilde\Phi\, u_R^g
                    - (Y_d)_{fg}\, \bar Q_L^f\, \Phi\, d_R^g
                    - (Y_e)_{fg}\, \bar \ell_L^f\, \Phi\, e_R^g
                    + \text{h.c.} ,
    \label{eq:theory-yukawa}
\end{equation}
where $\tilde\Phi_i = \varepsilon_{ij}\Phi^{*}_j$ is the conjugate doublet
and $f,g$ are flavour indices. The matrices $Y$ are the reason flavour
indices need first-class support: they are generic $3\times 3$ complex
matrices, not multiples of the identity, and diagonalising them is what
introduces quark mixing.

Fermions are Grassmann-valued, so reordering two of them costs a sign. This
is a trivial statement on paper and a persistent source of difficulty in an
implementation, because the symbolic engine underneath multiplies
commutatively. The convention adopted, and the one case where it must be
overridden, are discussed in \Cref{sec:internals-contraction-fermions}.

\subsection{Gauge fixing and Faddeev--Popov ghosts}
\label{sec:theory-gaugefixing}

Gauge invariance, having been imposed, immediately causes a problem: the
gauge field's kinetic operator is not invertible, because field
configurations related by a gauge transformation are physically identical and
the path integral counts each of them infinitely many times. There is no
propagator until one representative per gauge orbit is selected.

The Faddeev--Popov procedure~\cite{Faddeev:1967fc} performs that selection by
inserting a gauge-fixing condition, at the price of a Jacobian. In a linear
covariant gauge the condition is $\partial^\mu A^a_\mu = 0$, imposed softly
through
\begin{equation}
    \mathcal{L}_{\mathrm{gf}}
    = -\frac{1}{2\xi} \left(\partial^\mu A^a_\mu\right)
                      \left(\partial^\nu A^a_\nu\right) ,
    \label{eq:theory-gf}
\end{equation}
where the gauge parameter $\xi$ interpolates between the Landau gauge
($\xi\to0$) and the Feynman gauge ($\xi=1$). Physical results do not depend
on $\xi$, which makes it a useful internal consistency check.

The Jacobian is represented by the ghost Lagrangian
\begin{equation}
    \mathcal{L}_{\mathrm{gh}} = -\bar c^{\,a}\, \partial^\mu (D_\mu c)^a ,
    \qquad
    (D_\mu c)^a = \partial_\mu c^a + g f^{abc} A^b_\mu c^c ,
    \label{eq:theory-ghost}
\end{equation}
whose fields $c^a, \bar c^{\,a}$ are scalars carrying an adjoint index but
obeying Fermi statistics. They are not physical particles; they exist to
cancel the unphysical polarisations of the gauge boson in loops. For an
abelian group the ghost decouples, because $f^{abc}=0$ removes the
interaction term entirely.

\subsection{BRST symmetry}
\label{sec:theory-brst}

Fixing the gauge destroys gauge invariance, but not all of it. What survives
is a global, Grassmann-odd symmetry of the gauge-fixed
action~\cite{Becchi:1975nq}, generated by the BRST operator $s$:
\begin{align}
    s A^a_\mu &= (D_\mu c)^a , &
    s c^a     &= -\tfrac{g}{2} f^{abc} c^b c^c , \\
    s \bar c^{\,a} &= B^a , &
    s B^a     &= 0 ,
\end{align}
with $B^a$ an auxiliary field, and $s\phi_i = i g\, c^a (T^a)_{ij}\phi_j$ on
matter. Its defining property is nilpotency, $s^2 = 0$, and its practical use
here is that the entire gauge-fixing and ghost sector can be written as a
BRST variation,
\begin{equation}
    \mathcal{L}_{\mathrm{gf}} + \mathcal{L}_{\mathrm{gh}}
    = s\!\left( \bar c^{\,a}\, \partial^\mu A^a_\mu
                + \frac{\xi}{2}\, \bar c^{\,a} B^a \right) ,
\end{equation}
which makes the invariance of the full action manifest. Because $s$ is an
operator acting on the fields of a Lagrangian, it is a natural thing to
implement as a rewriting rule; \Cref{sec:feynpy-brst} does exactly that and
verifies $s^2=0$ symbolically.

\subsection{Spontaneous symmetry breaking and the physical basis}
\label{sec:theory-ssb}

The Standard Model is written in a basis that does not describe observable
particles. The electroweak gauge group $\mathrm{SU}(2)_L\times
\mathrm{U}(1)_Y$ has gauge fields $W^I_\mu$ and $B_\mu$, none of which is the
photon, and all matter fields are massless because chiral gauge invariance
forbids mass terms.

The resolution is that the Higgs doublet $\Phi$ has a potential
\begin{equation}
    V(\Phi) = \mu^2 \Phi^\dagger\Phi + \lambda (\Phi^\dagger\Phi)^2
\end{equation}
whose minimum for $\mu^2<0$ lies away from the origin. Choosing a vacuum
breaks the symmetry down to the electromagnetic $\mathrm{U}(1)_{\mathrm{em}}$
and, in the parametrisation used here, the doublet is expanded around that
vacuum as
\begin{equation}
    \Phi = \begin{pmatrix} G^+ \\
      \tfrac{1}{\sqrt2}\left(v + H + i G^0\right) \end{pmatrix} ,
    \label{eq:theory-vev}
\end{equation}
with $H$ the Higgs boson, $G^\pm, G^0$ the Goldstone bosons and $v$ the
vacuum expectation value. The gauge fields recombine into the physical mass
eigenstates,
\begin{equation}
    W^\pm_\mu = \frac{1}{\sqrt2}\left(W^1_\mu \mp i W^2_\mu\right),
    \qquad
    \begin{pmatrix} Z_\mu \\ A_\mu \end{pmatrix}
    =
    \begin{pmatrix} c_W & -s_W \\ s_W & c_W \end{pmatrix}
    \begin{pmatrix} W^3_\mu \\ B_\mu \end{pmatrix},
    \label{eq:theory-mixing}
\end{equation}
where $s_W, c_W$ are the sine and cosine of the weak mixing angle and satisfy
$s_W^2 + c_W^2 = 1$ --- an identity that reappears in
\Cref{sec:validation-sm}, because two implementations may distribute it
differently across the same vertex. Diagonalising the Yukawa matrices of
\Cref{eq:theory-yukawa} finally rotates the quark flavour basis and leaves
the CKM matrix in the charged current.

For this thesis the important observation is structural. Every step above is
a \emph{substitution}: replace $B_\mu$ by a combination of $Z_\mu$ and
$A_\mu$, replace one component of $\Phi$ by a constant plus $H$ plus $iG^0$,
replace $Q_L$ by a projector times a rotated field. A toolkit that can
perform such substitutions on an already-compiled Lagrangian can therefore
handle symmetry breaking without any dedicated symmetry-breaking code, which
is the design adopted in \Cref{sec:feynpy-transformations}.

\subsection{The background field method}
\label{sec:theory-bfm}

One further rewriting deserves mention because it is used by one of the
validation benchmarks. In the background field
method~\cite{Abbott:1980hw,Denner:1994xt}, each gauge field is split into a
classical background piece and a quantum fluctuation,
\begin{equation}
    A^a_\mu \longrightarrow \hat A^a_\mu + A^a_\mu ,
\end{equation}
and the gauge is fixed only for the quantum part, using a condition that is
covariant with respect to the background field. The resulting effective
action retains explicit gauge invariance in the background field, which
simplifies renormalisation considerably.

Structurally this is again a substitution, of exactly the same kind as
\Cref{eq:theory-mixing}, and the model \path{models/UnbrokenSM_BFM/}
implements it that way.

% ==============================================================
\section{Effective Field Theory and the SMEFT}
\label{sec:theory-smeft}
% ==============================================================

The most demanding benchmark of this thesis is not the Standard Model but an
effective field theory built on top of it. This section fixes the ideas and
the vocabulary that \Cref{ch:validation} later uses without further
explanation. A comprehensive review is~\cite{Brivio:2017vri}.

\subsection{The effective expansion}

Suppose new physics exists at some scale $\Lambda$ well above the energies
currently accessible. Its effects at low energy can then be described without
knowing what it is, by supplementing the Standard Model Lagrangian with all
operators that respect its symmetries and are built from its fields, ordered
by mass dimension:
\begin{equation}
    \mathcal{L}_{\mathrm{SMEFT}}
    = \mathcal{L}_{\mathrm{SM}}
    + \sum_i \frac{C^{(5)}_i}{\Lambda} \mathcal{O}^{(5)}_i
    + \sum_i \frac{C^{(6)}_i}{\Lambda^2} \mathcal{O}^{(6)}_i
    + \dots
    \label{eq:theory-smeft}
\end{equation}
Each operator comes with a \emph{Wilson coefficient} $C_i$: a number, or,
when the operator carries flavour indices, a tensor, encoding the unknown
short-distance physics. Higher-dimension operators are suppressed by more
powers of $\Lambda$, so the series can be truncated.

At dimension five there is only one gauge-invariant structure, the Weinberg
operator~\cite{Weinberg:1979sa}, which violates lepton number and generates
neutrino masses; it appears in the validation of
\Cref{sec:validation-ec-pinned}. At dimension six the count explodes, and
this is where the automation argument of \Cref{sec:intro-motivation} becomes
decisive.

\subsection{Bases, redundancy and the Green basis}
\label{sec:theory-green-basis}

Not every operator one can write is independent. Three relations reduce the
list. Integration by parts discards total derivatives. Fierz identities
relate different orderings of four-fermion structures. Most importantly, the
leading-order equations of motion can be used to eliminate operators, because
an operator proportional to the equations of motion can be removed by a field
redefinition and therefore does not affect $S$-matrix elements.

Applying all three yields a minimal, non-redundant set. The standard such set
at dimension six is the \emph{Warsaw basis}~\cite{Grzadkowski:2010es}, which
contains $59$ independent operators for one fermion generation, and $2499$
once flavour indices are counted. It is the basis in which most
phenomenological analyses are expressed. An earlier, redundant enumeration
was given in~\cite{Buchmuller:1985jz}.

For some purposes, however, one deliberately does \emph{not} reduce. A
\emph{Green basis}, also called an off-shell basis, retains the operators
that the equations of motion would have
removed~\cite{Gherardi:2020det,Chala:2021cgt}. This is redundant by
construction, and that is the point: off-shell quantities such as Green
functions, and the intermediate steps of matching and renormalisation
calculations~\cite{Jenkins:2013zja}, are not invariant under the field
redefinitions used to reach the Warsaw basis. Working in a Green basis first,
and reducing to a physical basis afterwards, keeps those intermediate steps
meaningful.

The model validated in \Cref{ch:validation} is of exactly this kind. It is a
baryon-number-preserving Green-basis SMEFT with $298$ active Wilson
coefficients spread over $32$ Lagrangian blocks --- larger than the Warsaw
basis precisely because it is redundant.

\subsection{Evanescent operators}
\label{sec:theory-evanescent}

One last complication is specific to the regularisation scheme. Calculations
are performed in $d = 4 - 2\epsilon$ dimensions, and Dirac algebra identities
that hold in exactly four dimensions --- such as the relations that allow a
product of two currents to be Fierz-reordered --- do not hold for general
$d$. Operator structures whose difference vanishes as $d\to4$, but not
before, are called \emph{evanescent}
operators~\cite{Buras:1998raa,Dugan:1990df}.

They cannot simply be discarded. Their matrix elements are of order
$\epsilon$, which combines with the $1/\epsilon$ poles of loop integrals to
leave a finite contribution, so a consistent scheme must carry them
explicitly. In the reference model they appear as a family of four-fermion
operators whose Wilson coefficients are named with the prefix
\texttt{alphaEc}, and they turn out to be exactly the operators where the two
implementations package charge conjugation differently. That is the subject
of \Cref{sec:validation-ec}.

% ==============================================================
\section{From the Lagrangian to Feynman Rules}
\label{sec:feynrules-algorithm}

In this section we want to explain but algorithm has been used to calculate Feynamn Rules given a lagrangian. If the reader lacks of kwoledge in QFT he is very welcomed to read a good textbook on the topic. To fully understand the algorithm the reader should be familiar with canonical quantization, the interaction picture, or Wick's theorem, we
assume this material here and use it only to fix notation.

In the previous section, we learnt about the role of field indices and Lagrangians. We can therefore move on to the central question: given a Lorentz-invariant and gauge-invariant Lagrangian, how do we obtain the associated Feynman rules? We wish to find an implementation that solves this in a systematic and mechanical way, so that it is suitable for implementation in Python or any programming language.
This implementation has already been provided by \textsc{FeynRules}~\cite{Christensen:2008py}, and the algorithm
described below reproduces, with some additional details, the
derivation given in Section~9 of that article.


\subsection{Mode Expansion and Canonical (Anti)Commutators}

Let's start by considering a free field $\phi^\alpha$, where the index $\alpha$ collects
all indices and quantum numbers carried by the field (Lorentz, Dirac,
and internal). In the interaction picture, we can $\phi^\alpha$ expand
the field in creation and annihilation operators,
\begin{equation}
    \phi^\alpha(x) = \int \! d\Pi(p) \,
    \Big( u^\alpha(p)\, a_\phi(p)\, e^{-ipx}
        + v^\alpha(p)\, a_{\bar\phi}^\dagger(p)\, e^{+ipx} \Big),
\end{equation}
where $u^\alpha(p)$, $v^\alpha(p)$ are the wavefunctions associated to
$\phi^\alpha$, $a_\phi(p)$ annihilates $\phi$, and
$a_{\bar\phi}^\dagger(p)$ creates a conjugated field
$\bar\phi$. Canonical quantization fixes the (anti)commutation
relations between fields and creation operators; schematically,
\begin{equation}\label{comm_rule}
    \big[\phi^\alpha(x),\, a_{\bar\phi}^{\dagger\,\beta}(p)\big]_{\pm}
    = \delta_{\alpha\beta} u^\alpha(p)\, e^{-ipx},
\end{equation}
with a commutator for bosonic fields and  anticommutator for
fermionic fields.


\subsection{The FeynRules Algorithm}

In order to understand the algorithm let's start by considering a general single term of a Lagrangian
\begin{equation}\label{eq:general-term}
    \mathcal L \supset g_{\alpha_1 \cdots \alpha_n}\,
    \partial \cdots \partial\, \phi^{\alpha_1} \cdots \phi^{\alpha_n},
\end{equation}
where $\alpha_i$ represents all the indices and quantum numbers of the field $\phi^{\alpha_i}$. $g_{\alpha_1\cdots\alpha_n}$ is a coupling constant (possibly
carrying its own indices), and $\partial \cdots \partial$ denotes some
number of derivatives acting on one or more of the fields. The
corresponding Feynman rule is obtained in three steps.

\begin{enumerate}
\item \textbf{Attach creation operators.} Multiply the term on the
right by one creation operator for each field, with all momenta taken
ingoing, and in the reverse order to that in which the (fermionic and
ghost) fields appear in the Lagrangian term. Sandwich the result
between vacuum states,
\begin{equation}
    g_{\alpha_1\cdots\alpha_n}\, \partial\cdots\partial\,
    \big\langle 0 \big|\,
    \phi^{\alpha_1}(x) \cdots \phi^{\alpha_n}(x)\,
    a^\dagger_{\bar\phi_n}(p_n) \cdots a^\dagger_{\bar\phi_1}(p_1)
    \,\big| 0 \big\rangle.
\end{equation}


\item \textbf{Contract creation operators to the left.} Now every creation operator is moved to the left, using the (anti)commutation rules of the fields given by equation \ref{comm_rule}. Note that moving $a^\dagger$ to the left of a boson produces a commutator term; moving it past a fermion produces an anticommutator term and a sign from reordering the
remaining operators. Concretely, each step replaces a pair
$\big(\phi^{\alpha_i}(x), a^\dagger_{\bar\phi_j}(p)\big)$ by the
contraction
\begin{equation}
    \delta_{\alpha_i\beta_j}\, u^{\alpha_i}(p)\, e^{-ipx},
\end{equation}

Once a creation operator reaches the leftmost position, it annihilates the vacuum $\langle 0|$ and is removed.

\item \textbf{Strip the overall factors.} After every creation
operator has been contracted away, the surviving vacuum overlap is
$\langle 0 | 0 \rangle = 1$. The remaining derivatives act on the
exponentials $e^{-ip_i x}$, pulling down the corresponding momenta.
Dropping the overall exponential $\exp\!\big(-i(p_1 + \cdots +
p_n)x\big)$ and the external wavefunctions
$u^{\alpha_i}(p_i)$, and multiplying the result by $i$, yields the
Feynman rule associated to the vertex.
\end{enumerate}

This procedure is nothing other than an automated application of
Wick's theorem to a single interaction term.

\subsection{Example: the QED Vertex}

Now we want to see how this algorithm applies to a Lagrangian. AS in ~\cite{Christensen:2008py} we illustrate the algorithm on the photon-lepton interaction term of
the QED Lagrangian,
\begin{equation}
    \mathcal L_{\bar\psi\psi A}
    = -e\, \bar\psi \gamma^\mu \psi\, A_\mu
    = -e\, \bar\psi_{s,f}\, \gamma^\mu_{ss'}\, \psi_{s',f}\, A_\mu,
\end{equation}
where $e$ is the electromagnetic coupling, $\psi_f$ the lepton field
of flavour $f \in \{e,\mu,\tau\}$, and $A_\mu$ the photon field.
Multiplying by the creation operators $a^\dagger_{\bar\psi}$,
$a^\dagger_\psi$, $a^\dagger_A$ and sandwiching between vacuum states gives
\begin{equation}
    -e\, \big\langle 0 \big|\,
    \bar\psi \gamma^\mu \psi\, A_\mu\,
    a^\dagger_{\bar\psi}\, a^\dagger_\psi\, a^\dagger_A
    \,\big| 0 \big\rangle.
\end{equation}
Moving each creation operator to the left using the canonical
(anti)commutators,
\begin{align}
    \big\{ \psi_{s,f}(x), a^\dagger_{\bar\psi_{s',f'}}(p) \big\}
        &= \delta_{ss'}\delta_{ff'}\, u_{s,f}\, e^{-ipx}, \\
    \big\{ \bar\psi_{s,f}(x), a^\dagger_{\psi_{s',f'}}(p) \big\}
        &= \delta_{ss'}\delta_{ff'}\, \bar u_{s,f}\, e^{-ipx}, \\
    \big[ A_\mu(x), a^\dagger_{A_\nu}(p) \big]
        &= \delta_\mu^\nu\, \epsilon_\mu\, e^{-ipx},
\end{align}
and annihilating the vacuum, yields
\begin{equation}
    -e\, \delta_{ff'}\, \gamma^\mu_{ss'}\,
    \bar u_{sf}(p_1)\, u_{s'f'}(p_2)\, \epsilon_\mu(p_3)\,
    e^{-i(p_1+p_2+p_3)x}.
\end{equation}
Dropping the wavefunctions and the overall exponential, and
multiplying by $i$, gives the familiar QED vertex,
\begin{equation}
    -ie\, \delta_{ff'}\, \gamma^\mu_{ss'}.
\end{equation}

% ==============================================================
\section{Conventions}
\label{sec:theory-conventions}
% ==============================================================

Almost every disagreement between two independently written derivations of a
Feynman rule turns out, on inspection, to be a disagreement about a
convention rather than about physics. \Cref{ch:validation} is largely an
exercise in identifying which ones. It is therefore worth collecting in one
place the choices that the implementation makes and enforces everywhere;
\Cref{tab:theory-conventions} is referred back to from
\Cref{ch:internals,ch:validation}.

\begin{table}[htbp]
\centering
\begin{tabular}{@{}p{4.5cm}p{8.5cm}@{}}
\toprule
\textbf{Quantity} & \textbf{Convention} \\
\midrule
Metric signature
    & $g_{\mu\nu} = \mathrm{diag}(+,-,-,-)$ \\
Covariant derivative (matter)
    & $D_\mu = \partial_\mu - i g A^a_\mu T_a$, \Cref{eq:theory-covd} \\
Field strength
    & $F^a_{\mu\nu} = \partial_\mu A^a_\nu - \partial_\nu A^a_\mu
       + g f^{abc} A^b_\mu A^c_\nu$, \Cref{eq:theory-fs} \\
Adjoint covariant derivative
    & $(D_\mu X)^a = \partial_\mu X^a + g f^{abc} A^b_\mu X^c$ \\
Generator normalisation
    & $[T_a, T_b] = i f_{ab}{}^{c} T_c$, generators hermitian \\
Fourier convention
    & all momenta incoming; $\partial_\mu \to -i p_\mu$ \\
Overall vertex factor
    & every extracted rule is multiplied by $+i$ \\
Chiral projectors
    & $P_{L,R} = (1 \mp \gamma^5)/2$, \Cref{eq:theory-projectors} \\
Charge conjugation
    & $C^{T} = -C$ \\
Fermion ordering
    & the order in which fermion fields are written in a Lagrangian term is
      significant and is preserved throughout the pipeline \\
\bottomrule
\end{tabular}
\caption{The conventions fixed once and used throughout the thesis and the
implementation. The signs in the first four rows are choices, not results;
what matters is that they are applied uniformly, since a single inconsistency
would produce vertices that are individually plausible and collectively
wrong.}
\label{tab:theory-conventions}
\end{table}

With the theory, the algorithm and the conventions in place, we can turn to
the toolkit itself.
